\documentclass[aspectratio=169,10pt]{beamer}

\usetheme{metropolis}
\usefonttheme[onlymath]{serif}
\setbeamertemplate{navigation symbols}{}

\usepackage[UTF8]{ctex}
\usepackage{amsmath, amssymb, amsfonts, bm}
\usepackage{graphicx}
\usepackage{booktabs}
\usepackage{hyperref}

\def\dif{\mathinner{\mathrm{d}}\hphantom{\mskip-\thinmuskip}}

\title{基于随机特征法的高维问题求解研究}
\subtitle{本科毕业论文开题答辩}
\author{学生姓名: 刘行}
\institute{中国科学技术大学大学\quad 少年班学院}
\date{\today}

\begin{document}
    \begin{frame}
        \titlepage
    \end{frame}

    \begin{frame}{报告提纲}
        \tableofcontents
    \end{frame}

    \section{研究背景与意义}
    \begin{frame}{研究背景}
        \begin{itemize}
            \item 高维偏微分方程在金融数学, 量子物理和控制问题中普遍存在.
            \item 传统数值方法 (有限元法, 有限差分法) 在高维情形下计算复杂度指数增长, 出现 ``维度灾难''.
            \item 随机特征法 (Random Feature Method, RFM)与深度学习相结合, 为高维问题提供了新的近似手段.
        \end{itemize}
    \end{frame}

    \begin{frame}{研究意义}
        \begin{itemize}
            \item 探索基于随机特征的高维数值求解方法, 可突破传统网格法的限制.
            \item 为高维抛物型方程提供新的数值框架.
            \item 对比分析 Deep BSDE 方法与 RFM 方法的性能差异, 为后续研究提供参考.
        \end{itemize}
    \end{frame}

    \section{研究目标与主要内容}
    \begin{frame}{研究目标}
        \begin{enumerate}
            \item 构建随机特征法 (RFM) 求解高维 PDE 的通用框架.
            \item 实现半线性抛物型方程的 Deep BSDE 求解器.
            \item 设计线性抛物型方程的 RFM 求解器并进行实验验证.
        \end{enumerate}
    \end{frame}

    \begin{frame}{主要研究内容}
        \begin{columns}[T]
            \column{0.5\textwidth}
            \begin{block}{理论部分}
                \begin{itemize}
                    \item PDE 与 BSDE 的等价关系
                    \item 随机特征法的数学基础
                \end{itemize}
            \end{block}

            \column{0.5\textwidth}
            \begin{block}{实验部分}
                \begin{itemize}
                    \item Deep BSDE 求解器的实现与优化
                    \item 线性 PDE 的 RFM 实验框架
                    \item 加入非线性求解器以求解半线性 PDE
                    \item 对比两种方法的性能
                \end{itemize}
            \end{block}
        \end{columns}
    \end{frame}

    \section{相关方法简介}
    \begin{frame}{PDE 与 BSDE 的关系}
        考虑半线性抛物型方程:
        \begin{equation*}
            \frac{\partial u}{\partial t} + \frac{1}{2} \mathrm{Tr}\left(\sigma\sigma^{\top}\,\mathrm{Hess}_x u\right) + \langle\nabla u, \mu\rangle + f\left(t, x, u, \sigma^{\top} \nabla u\right)=0, \quad u\left(T,x\right)=g\left(x\right)
        \end{equation*}
        对应的后向随机微分方程 (BSDE) 为:
        \begin{equation*}
            \left\{\begin{array}{l}
                X_{t} = \xi + \int_{0}^{t} \mu\left(s, X_{s}\right) \dif s + \int_{0}^{t} \sigma\left(s, X_{s}\right) \dif W_{s} \\
                Y_{t} = g\left(X_{T}\right) + \int_{t}^{T} f\left(s, X_{s}, Y_{s}, Z_{s}\right) \dif s - \int_{t}^{T} \left(Z_{s}\right)^{\top} \dif W_{s}
            \end{array}\right.
        \end{equation*}
        其中 $Y_{t} = u\left(t, X_{t}\right), Z_{t} = \left[\sigma\left(t, X_{t}\right)\right]^{\top} \left(\nabla_{x} u \right)\left(t, X_{t}\right)$, 则随机过程 $\left(X_{t}, Y_{t}, Z_{t}\right) \in \mathbb{R}^{d} \times \mathbb{R} \times \mathbb{R}^d, t \in \left[0, T\right]$.

        为简化模型, 通常取 $\mu = 0$, $\sigma = I_d$ 或 $\sigma = \sqrt{2} I_d$ 等.
    \end{frame}

    \begin{frame}{随机特征法 (RFM) 简介}
        \begin{itemize}
            \item 基本思想: 用随机基函数近似高维空间中的未知函数.
            \item 表示形式:
            \begin{equation*}
                u\left(x\right) \approx \sum_{i=1}^{m} a_{i} \, \phi\left(\omega_{i}^{\top} x + b_{i}\right)
            \end{equation*}
            其中参数 $\omega_{i}, b_{i}$ 随机生成，系数 $a_{i}$ 通过最小二乘拟合.
            \item 优点: 避免高维网格划分, 训练速度快, 可扩展至非线性 PDE.
        \end{itemize}
    \end{frame}

    \section{已完成工作}
    \begin{frame}{Deep BSDE 求解器}
        \begin{itemize}
            \item 已完成半线性抛物型方程的 Deep BSDE 求解器实现, 使用时间离散 + 神经网络近似 $Z_{t}$.
            \item 损失函数来源于终端条件误差: $\mathcal{L} = \mathbb{E}\left\lvert Y_{T} - g\left(X_{T}\right)\right\rvert^{2}$
        \end{itemize}

        \vspace{1em}
        \begin{figure}[htbp]
            \centering
            \includegraphics[width=0.5\textwidth]{figures/bsde_loss_curve.png}
        \end{figure}
    \end{frame}

    \begin{frame}{RFM 求解器 (进行中)}
        \begin{itemize}
            \item 已写出迭代过程, 但是在将正向迭代拼接为矩阵以实现最小二乘反向求解时出错, 正在重新整理思路.
            \item 当前实现:
                \begin{itemize}
                    \item 取随机特征函数 $\boldsymbol{\phi}\left(x, t\right) = \alpha\tanh\left(A \boldsymbol{x} + bt + c\right)$
                    \item 随机采样特征参数 $(\omega_{i}, b_{i})$
                    \item 构造随机特征矩阵 $\Phi$
                    \item 通过最小二乘求解系数
                \end{itemize}
        \end{itemize}
        \begin{equation*}
            \begin{array}{l}
                \dif w^{k}, x^{k} = \mathrm{select()} \\
                \mathrm{for} \; n = 0, 1, \ldots, N-1: \\
                \quad \quad t_{n} = n \Delta t, \boldsymbol{x}_{n}^{k} = x^{k}\left[:,n\right], \dif \boldsymbol{w}_{n}^{k} = \dif w^{k}\left[:,n\right] \\
                \quad \quad \boldsymbol{z}_{n}^{k} = \boldsymbol{\phi}\left(\boldsymbol{x}_{n}^{k}, t_{n}\right) \\
                \quad \quad y_{n+1}^{k} = y_{n}^{k} - \Delta t \cdot f\left(t_{n}, \boldsymbol{x}_{n}^{k}, y_{n}^{k}, \boldsymbol{z}_{n}^{k}\right) + \left\langle \dif \boldsymbol{w}_{n}^{k}, \boldsymbol{z}_{n}^{k} \right\rangle \\
            \end{array}
        \end{equation*}
    \end{frame}

    \section{后续工作与计划}
    \begin{frame}{后续工作计划}
        \begin{enumerate}
            \item 完成线性 PDE 的 RFM 数值实验与误差分析.
            \item 推广至非线性方程, 比较与 Deep BSDE 方法.
            \item 优化采样与正则化策略, 提高稳定性.
            \item 撰写论文并准备毕业答辩.
        \end{enumerate}
    \end{frame}

    \section{总结与展望}
    \begin{frame}{总结与展望}
        \begin{itemize}
            \item 已构建 Deep BSDE 求解器并验证可行性;
            \item RFM 求解器框架理论部分基本完成, 后续将完善实验;
            \item 研究有望为高维 PDE 数值解提供新的高效思路.
        \end{itemize}
        \vfill
        \centering
        {\large 谢谢!}
    \end{frame}
\end{document}
